\documentclass[11pt]{article}
\usepackage[utf8]{inputenc}
\usepackage[margin=1in]{geometry}
\usepackage{amsmath,amssymb}
\usepackage{booktabs}
\usepackage{hyperref}
\usepackage{natbib}
\usepackage{setspace}
\usepackage{float}
\usepackage{enumitem}
\onehalfspacing

\hypersetup{colorlinks=true,linkcolor=blue!70!black,citecolor=blue!70!black,urlcolor=blue!70!black}

\title{\textbf{Regime-Aware Factor Backtest:\\Point-in-Time Equity Strategy with\\HMM Regime Detection and Fama-French Attribution}}
\author{Tanishk Yadav\\[2pt]\textit{MS Financial Engineering}\\NYU Tandon School of Engineering\\New York, NY, USA\\[2pt]ORCID: 0009-0006-2382-9411}
\date{}

\begin{document}
\maketitle

\begin{abstract}
\noindent
We present a deliberately bias-free quantitative equity backtest that integrates
walk-forward Hidden Markov Model (HMM) regime detection with SEC filing-gated
fundamental data and Fama-French 5-Factor (FF5) performance attribution. Fundamental
data is sourced through ChronoFund with zero look-ahead via SEC
\texttt{acceptance\_datetime} gating; regimes are fit only on trailing data; held-name
delistings are penalized; and all returns are net of a 2-and-15 fee structure (2\%
management, 15\% performance with high-water mark). Our central result is an honest
one: \emph{alpha decay}. Over 2008-2026 the strategy's FF5 alpha is statistically
\emph{insignificant} ($2.42\%$/yr, $t = 0.91$, $p = 0.365$). Splitting the sample
reveals why: alpha was significant in the first half (2008-2017: $8.37\%$/yr,
$t = 2.73$, $p = 0.006$) and vanished in the second (2017-2026: $-3.44\%$/yr,
$p = 0.435$). Net of fees the strategy still edges SPY on CAGR ($12.7\%$ vs $11.4\%$),
Sharpe ($0.61$ vs $0.57$), and maximum drawdown ($-39.7\%$ vs $-51.5\%$), but FF5
attribution shows this edge is market beta ($0.85$) and a small-cap tilt
($\text{SMB} = 0.18$), not residual alpha. The contribution is the reproducible,
bias-free framework and the disciplined out-of-sample attribution, not the headline
return.

\medskip
\noindent\textbf{Keywords:} quantitative equity, regime detection, Fama-French, point-in-time data, alpha attribution, SEC EDGAR
\end{abstract}

\section{Introduction}

This paper describes a systematic equity strategy that synthesizes three research threads: (1)~regime-aware portfolio construction using Hidden Markov Models, (2)~point-in-time fundamental data engineering via SEC EDGAR, and (3)~rigorous performance attribution using the Fama-French 5-Factor model.

The strategy's defining characteristic is its commitment to bias-free research. Every fundamental data point is gated by the SEC's \texttt{acceptance\_datetime} - the precise second the filing was accepted - through the ChronoFund data engine. The HMM regime detector uses walk-forward estimation with expanding windows. Transaction costs, management fees, and performance fees are modeled explicitly.

\section{Data and Infrastructure}

\subsection{Fundamental Data}
Company fundamentals (income statements, balance sheets, cash flow statements, and derived metrics) are sourced from SEC EDGAR via the ChronoFund engine. The point-in-time constraint ensures that for any backtest date $t$, only filings with \texttt{acceptance\_datetime} $\leq t$ are available. Historical data includes pre-computed Parquet files: \texttt{company\_master}, \texttt{filings}, \texttt{statements\_income}, \texttt{statements\_balance}, \texttt{statements\_cashflow}, and \texttt{derived\_metrics}.

\subsection{Price Data}
Daily adjusted close prices for all universe constituents and the SPY benchmark.

\subsection{Factor Data}
Daily Fama-French 5-Factor data (Mkt-RF, SMB, HML, RMW, CMA) stored as a Parquet file for efficient access during attribution.

\section{Methodology}

\subsection{Regime Detection}
A 4-state Gaussian HMM classifies market conditions into bull, normal, bear, and crisis states. The model is re-estimated at each rebalancing date using only historical data, with regime assignment based on filtered (not smoothed) probabilities.

\subsection{Ranking System}
The adaptive ranking engine (\texttt{rank\_system\_v2.py}, 317 lines) scores universe constituents using fundamental factors derived from the ChronoFund snapshot. The ranking weights are regime-conditional: growth factors receive higher weight in bull regimes, while quality and value factors are emphasized during bear and crisis regimes.

\subsection{Portfolio Construction}
The strategy holds a concentrated portfolio of top-ranked stocks, with position sizes determined by the ranking score and regime-dependent risk constraints. During crisis regimes, equity exposure is reduced and replaced with treasury positions.

\subsection{Options Hedging}
An optional put-spread overlay (\texttt{options\_hedge.py}, 795 lines) provides tail risk protection during elevated-stress regimes. The hedging module logs all hedge transactions and their P\&L contribution separately.

\subsection{Fee Structure}
\begin{itemize}[nosep]
    \item \textbf{Management fee}: 2\% per annum, charged monthly on AUM
    \item \textbf{Performance fee}: 15\% of profits above high-water mark, charged quarterly
    \item \textbf{Transaction costs}: Modeled per-trade based on estimated market impact
\end{itemize}

\section{Results}

\subsection{Performance Summary}

\begin{table}[H]
\centering
\begin{tabular}{lccc}
\toprule
\textbf{Metric} & \textbf{Strategy} & \textbf{SPY} & \textbf{$\Delta$} \\
\midrule
Total Return & \textbf{763.0\%} & 592.4\% & +170.6\% \\
Annual Return & \textbf{12.73\%} & 11.36\% & +1.37\% \\
Sharpe Ratio & \textbf{0.61} & 0.57 & +0.04 \\
Max Drawdown & \textbf{$-$39.74\%} & $-$51.48\% & +11.74\% \\
Volatility & 20.82\% & 19.89\% & +0.93\% \\
\bottomrule
\end{tabular}
\caption{Net-of-fee performance, 2008-2026. Fees: 2\% management + 15\% performance (HWM).}
\end{table}

\subsection{Fee Breakdown}
Total costs over the 18-year period on \$500M initial capital: \$60.0M transaction costs, \$1,003.8M management fees, \$1,007.9M performance fees. Total: \$2,071.7M. The strategy's gross performance is substantially higher; the fee drag is intentionally heavy to present conservative results.

\subsection{Fama-French 5-Factor Attribution}

\begin{table}[H]
\centering
\begin{tabular}{lcccc}
\toprule
& \textbf{Full Period} & & \textbf{First Half} & \textbf{Second Half} \\
& (2008-2026) & & (2008-2017) & (2017-2026) \\
\midrule
Alpha (ann.) & 2.42\% & & \textbf{8.37\%} & $-$3.44\% \\
$t$-stat & 0.91 & & \textbf{2.73} & $-$0.78 \\
$p$-value & 0.365 & & \textbf{0.006} & 0.435 \\
Mkt-RF $\beta$ & 0.85 & & 0.84 & 0.86 \\
$R^2$ & 0.708 & & 0.791 & 0.637 \\
\bottomrule
\end{tabular}
\caption{Fama-French 5-Factor regression. Alpha is statistically significant ($p = 0.006$) in the first half but decays in the second half.}
\end{table}

The first-half alpha of 8.37\% is statistically significant at the 1\% level. The second-half alpha is negative and insignificant, indicating that the fundamental factors captured by the ranking system have lost predictive power in the more recent period. The defensive beta (0.85) is consistent across both halves, confirming that the regime-dependent risk management is stable.

\paragraph{Adaptive response to decay.} The framework includes optional adaptive overlays that respond to the decay rather than merely diagnosing it. Relative to the static baseline (12.73\% ann., Sharpe 0.61, max drawdown $-$39.7\%), an alpha-fade overlay that blends toward the benchmark when rolling alpha turns negative raises annual return to 13.68\% and Sharpe to 0.66 while roughly halving the second-half alpha drag ($-$3.3\% to $-$1.6\%); a volatility-targeting overlay cuts maximum drawdown to $-$28.8\% at a higher Sharpe (0.64) and nearly eliminates the second-half drag. Importantly, even with these overlays the full-period FF5 alpha remains statistically insignificant: the overlays mitigate the decay's drag, they do not restore alpha.

\subsection{Regime Distribution}
Over the full period: 2,004 bull days (44.1\%), 1,232 normal days (27.1\%), 819 bear days (18.0\%), 476 crisis days (10.5\%). The 230 trades in 18 years imply an average holding period of approximately one month.

\section{Limitations}

\begin{itemize}[nosep]
    \item \textbf{Alpha decay}: The strategy's alpha is concentrated in the first half. This may reflect factor crowding, changing market microstructure, or the specific factor weights becoming stale.
    \item \textbf{Fixed factor weights}: The ranking system uses regime-conditional but otherwise fixed factor weights. Adaptive weight learning could improve the second-half performance.
    \item \textbf{Sector-level execution}: Assumes perfect liquidity in ETFs, which mostly holds but degrades for smaller sectors during stress.
    \item \textbf{Gaussian HMM}: Inherits all limitations of the Gaussian emission assumption from the underlying regime detection model.
    \item \textbf{Survivorship bias}: While ChronoFund tracks delisted companies, the current universe construction does not fully account for index reconstitution effects.
\end{itemize}

\section{Conclusion}

The Regime-Aware Factor Backtest's primary contribution is methodological: a
reproducible, bias-free framework (point-in-time SEC-gated fundamentals, walk-forward
regimes, delisting penalties, realistic 2-and-15 fees) for honest factor research.
Its empirical result is a cautionary one: a leverage-averse, regime-aware fundamental
strategy that produced significant FF5 alpha in 2008-2017 saw that alpha fully decay
post-2017, leaving full-period alpha statistically insignificant. The strategy's
net-of-fee outperformance of SPY is real but is attributable to market beta and a
small-cap tilt rather than skill. This is exactly the kind of finding that
out-of-sample, fee-aware, look-ahead-free evaluation is designed to surface - and
that less disciplined backtests routinely hide. The 3,582-line codebase and its
integration with the ChronoFund data engine make the result fully reproducible.

\medskip
\noindent\textbf{Code Availability:} \url{https://github.com/tanishhky/Regime-Aware-Factor-Backtest}

\begin{thebibliography}{5}
\bibitem{fama2015} E.~F. Fama and K.~R. French, ``A five-factor asset pricing model,'' \textit{Journal of Financial Economics}, vol.~116, no.~1, pp.~1-22, 2015.
\bibitem{hamilton1989} J.~D. Hamilton, ``A new approach to the economic analysis of nonstationary time series,'' \textit{Econometrica}, vol.~57, no.~2, pp.~357-384, 1989.
\bibitem{novy2012} R.~Novy-Marx, ``The other side of value: The gross profitability premium,'' \textit{Journal of Financial Economics}, vol.~108, no.~1, pp.~1-28, 2013.
\end{thebibliography}

\end{document}
